
% !TEX program = lualatex
\documentclass[a4paper,12pt]{bxjsbook}

% =========================================================
% Packages
% =========================================================

% --- Japanese (LuaLaTeX) ---
\usepackage{luatexja}

% --- math / layout ---
\usepackage{amssymb}
\usepackage{amsthm}
\usepackage{mathtools}
\usepackage{geometry}
\usepackage{mathpartir}
\geometry{margin=25mm}

% --- lists / links / figures ---
\usepackage{enumitem}
\usepackage{xurl}
\usepackage{hyperref}
\usepackage{tikz}
\usetikzlibrary{arrows.meta, positioning, calc, fit, backgrounds, decorations.markings}

\hypersetup{
  colorlinks=true,
  linkcolor=blue,
  urlcolor=blue,
  citecolor=blue
}

% =========================================================
% length definition
% =========================================================
\newlength{\customindent}
\setlength{\customindent}{2em}

% =========================================================
% macro definition
% =========================================================
\newcommand{\BodyListFixBegin}{%
  \begingroup
  \setlength{\rightskip}{0pt}%
  \setlength{\parindent}{0pt}%
  % \setlength{\displaywidth}{\linewidth}%
}
\newcommand{\BodyListFixEnd}{\ifhmode\par\fi\endgroup}

\newenvironment{BodyListFix}{\BodyListFixBegin}{\BodyListFixEnd}

% =========================================================
% environment definition
% =========================================================

% Main_Sentense styles
\newenvironment{body}{
  \par
  \begingroup
  \setlength{\leftskip}{\customindent}
  \setlength{\rightskip}{0pt}
  \setlength{\parindent}{1em}
  \setlength{\displayindent}{\customindent}
  \setlength{\displaywidth}{\dimexpr\linewidth-\customindent\relax}
}{
  \par
  \endgroup
}

% description in Main_Sentense styles
\newlist{bodydescription}{description}{1}
\setlist[bodydescription]{
  labelwidth=0pt,
  labelsep=0pt,
  itemindent=0pt,
  style=nextline,
  font=\bfseries
}

% itemize in Main_Sentense styles
\newlist{bodyitemize}{itemize}{3}
\setlist[bodyitemize]{
  leftmargin=3em,
  label=\textbullet,
  before=\begin{BodyListFix},
  after=\end{BodyListFix}
}
\setlist[bodyitemize,2]{leftmargin=3.5em}
\setlist[bodyitemize,3]{leftmargin=4em}

% enumerate in Main_Sentense styles
\newlist{bodyenumerate}{enumerate}{3}
\setlist[bodyenumerate]{
  leftmargin=4em,
  label=(\arabic*),
  before=\begin{BodyListFix},
  after=\end{BodyListFix}
}
\setlist[bodyenumerate,2]{leftmargin=4.5em, label=(\alph*)}
\setlist[bodyenumerate,3]{leftmargin=5em, label=(\roman*)}

% Theorem styles
\newtheoremstyle{theorem}
  {\baselineskip}
  {\baselineskip}
  {
    \normalfont
    \setlength{\leftskip}{\customindent}
    \setlength{\parindent}{0pt}
    \setlength{\rightskip}{0pt}
    \setlength{\displayindent}{\customindent}
    \setlength{\displaywidth}{\dimexpr\linewidth-\customindent\relax}
  }
  {0pt}
  {\bfseries}
  {.}
  {\newline}
  {
    \hskip-\customindent
    \thmname{#1}\thmnumber{ #2}\thmnote{（#3）}
  }

\theoremstyle{theorem}
\newtheorem{theorem}{定理}[chapter]
\newtheorem{lemma}[theorem]{補題}
\newtheorem{proposition}[theorem]{命題}
\newtheorem{corollary}[theorem]{系}
\newtheorem{definition}[theorem]{定義}
\newtheorem{example}[theorem]{例}
\newtheorem{remark}[theorem]{注意}
\newtheorem{abbreviation}[theorem]{略記}

% description in Theorem styles
\newlist{theoremdescription}{description}{3}
\setlist[theoremdescription]{
  leftmargin=5em,
  label=\textbullet,
  before=\begin{BodyListFix},
  after=\end{BodyListFix}
}
\setlist[theoremdescription,2]{leftmargin=5.5em}
\setlist[theoremdescription,3]{leftmargin=6em}

% itemize in Theorem styles
\newlist{theoremitemize}{itemize}{3}
\setlist[theoremitemize]{
  leftmargin=5em,
  label=\textbullet,
  before=\begin{BodyListFix},
  after=\end{BodyListFix}
}
\setlist[theoremitemize,2]{leftmargin=5.5em}
\setlist[theoremitemize,3]{leftmargin=6em}

% enumerate in Theorem styles
\newlist{theoremenumerate}{enumerate}{3}
\setlist[theoremenumerate]{
  leftmargin=6em,
  label=(\arabic*),
  before=\begin{BodyListFix},
  after=\end{BodyListFix}
}
\setlist[theoremenumerate,2]{leftmargin=6.5em, label=(\alph*)}
\setlist[theoremenumerate,3]{leftmargin=7em, label=(\roman*)}

% =========================================================
% Proof Environment Setting
% =========================================================
\makeatletter
\renewenvironment{proof}[1][\proofname]{%
  \par\pushQED{\qed}%
  \normalfont
  \topsep6\p@\@plus6\p@\relax
  \list{}{%
    \leftmargin=\customindent
    \rightmargin=0pt
    \labelwidth=0pt
    \labelsep=0pt
    \itemindent=0pt
    \listparindent=0pt
    \parsep=0pt
  }%
  \item[\itshape #1\@addpunct{.}\hspace{0.5em}]%
  \setlength{\displayindent}{\customindent}%
  \setlength{\displaywidth}{\dimexpr\linewidth-\customindent\relax}%
  \ignorespaces
}{%
  \popQED\endlist\@endpefalse
}
\makeatother

% =========================================================
% Convenience commands
% =========================================================
\newcommand{\addchaptertotoc}[1]{%
  \chapter*{#1}%
  \addcontentsline{toc}{chapter}{#1}%
}
\newcommand{\dotminus}{\mathbin{\dot{-}}}

\DeclareMathOperator{\AffHull}{AffHull}
\DeclareMathOperator{\Span}{Span}
\DeclareMathOperator{\AffSpan}{AffSpan}
\DeclareMathOperator{\AffComb}{AffComb}
\DeclareMathOperator{\ConvHull}{ConvHull}
\DeclareMathOperator{\Face}{Face}
\DeclareMathOperator{\Facet}{Facet}



% =========================================================
% Metadata
% =========================================================
\title{チェバの定理の一般化と\\ Lean4/mathlib4への形式化}
\author{
秋田 隼\\
早稲田大学\\
基幹理工学部 数学科B4\\
1W221002-0
}

\date{$2026$/$1$/$18$}

% =========================================================
% Document
% =========================================================
\begin{document}

    \maketitle

    \frontmatter
    \addchaptertotoc{要旨}

    \begin{body}
        \indent 本研究は、古典的チェバの定理を出発点として、チェバの定理の一般化を定式化し、関連する補題のLean4/mathlib4上での形式化を目標とする。

        \vskip\baselineskip

        \noindent 主な貢献は以下である。

        \begin{bodyitemize}[leftmargin=3.5em]
            \item チェバの定理の一般化の定式化の紹介。
            \item 既存のmathlib4の構造を調査し、再利用可能部分と不足部分を切り分ける設計指針。
            \item 不足する補題・定義をモジュール化する方針を整理し、逆方向定理の機械検証に必要な論点を明示する。
        \end{bodyitemize}

        \vskip\baselineskip
        \noindent 実装はGitHubリポジトリで公開している：
        
        \url{https://github.com/Aj1905/LeanResearch.git}
    \end{body}

    \tableofcontents

    \mainmatter

    % =========================================================
    \chapter{序論}
    % =========================================================

    \section{研究背景と動機}
    \begin{body}
        \indent 本研究は、mathlib4に存在する順方向チェバの定理の形式化を基盤として、その定理の逆をLean4上で構成し、mathlib4への追加を目指す。mathlib4は、Lean Prover Communityにより開発される大規模数学ライブラリ\cite{LeanProverCommunity}であり、分野により整備状況には濃淡がある。
        特に幾何学は、形式化が代数的手法に比べて技術的に困難であることや、図形的直観の形式言語翻訳の複雑さが起因するためか発展が停滞気味である。こうした開発目標を可視化すべく、数学において重要でありながら未だ形式化されていない定理が「missing 100 theorems」としてリストアップされている。

        \indent チェバの定理は、図形の共点性を代数条件へ翻訳する基本的な道具であり、多数の幾何学的議論の分岐点として機能する。この種の定理を一般化してライブラリ化することの価値は、定理本体に加えて単体・面・アフィン結合・重心座標などの重要な周辺概念を再利用可能な形で整備できる点にある。よって今後の幾何定理を形式化する際の基盤を構築できるという点で大きな意義を持つと考える。

        \indent さらに、形式化は教科書的な証明に潜む、ゼロ除算回避、図形が退化する場合の排除、必要な代数構造などの暗黙の仮定を明示し、どの仮定のもとでどの結論が成立するかを定理の型として確定する。このことは、一般化の汎用性を証明の文章ではなくLeanの型と依存関係によって機械的に検証可能な形で提示できる点で、数学的にも工学的にも意義がある。近年もチェバの定理の一般化は研究されており、本研究はその形式的基盤をmathlib4に提供する試みとして位置づけられる。
    \end{body}

    \section{本稿の構成}
    \begin{body}
        \indent 第$2$章でLean4について解説し、第$3$章で一般化とは何か基礎論的な視点から眺める。第$4$章でチェバの定理とアフィン幾何の関連を議論し、第$5$章で古典的チェバの定理およびその逆がどのように一般化され得るかを考える。第$6$章で実際に行った実装の詳細を見る。

        \indent 本稿では暗黙に、自然演繹を推論規則として採用した古典一階述語論理のもとで議論を行なっている。集合論に言及する時、その概念の定義はKunenに従う。\cite{Kunen}

    \end{body}


    % =========================================================
    \chapter{Lean4による形式化}
    % =========================================================
    \section{形式言語の概要}
    \begin{body}
        \indent 自然言語の主張はさまざまな点で曖昧さを含む。定義が省略されていたり、暗黙の前提があったり、同じ記号が別の意味を持っていたりすることで本来したいはずの主張が伝わらないことや別の意味で解釈されてしまうこともある。

        \indent これを解決するために用いられるのが\emph{形式言語}と呼ばれる定義が厳密に定まった言語であり、自然言語の文を形式言語に翻訳する作業は\emph{形式化}と呼ばれる。一度形式化された文章は、読み手がその体系の定義を正確に把握している限り、いつ誰がどこで読んでも意味が一意に定まる。これはコンピュータが文章を読む場合も例外ではない。この形式言語の文章をどう分析しどう推論するかも含めて厳密に定義したものを形式体系と呼ぶ。

        \indent 一方、数学とは過去の数学者の成果の積み重ねからなる学問である。数学的主張は、他者に共有され、検証され、再利用される形で後世に伝え残されてはじめて真価を発揮する。このとき、時代や場所、読み手によって意味や解釈が揺れることは望ましくない。いついかなる時も同じ主張をすることが求められる。まさに形式体系の特徴が最大限活かされる分野の一つである。現代数学は素朴集合論を一階述語論理で公理化した形式体系（主にZFやZFC）を基礎として用いて構築されることが多い。

        \indent しかし一度考え直してみると、数学を記述する形式体系は必ずしもZFCに限る必要はない。ZFCは確かに人間が数学を理解する上で相性のいい形式体系の一種かもしれない。形式体系の中では比較的人間にとって理解しやすく、また歴史的に標準化されているため慣れてもいる。しかし本来、\emph{言葉は受取り手と目的に最適化することで意味伝達という役割を最大限発揮する}。伝えたい意味内容に対し、受取り手と目的を明確にして最適な形式体系を選択することが重要となる。こうしてZFCの代替案の一つとして挙げられるのが型理論である。

        \indent 簡潔に言えば型理論とは、命題を型、証明をその型の項として表現するという特徴を持った形式体系である。集合論において全ての数学的対象が集合であるように、型理論において全ての数学的対象は型を持つ。詳細な説明は\hyperref[appendix:type-theory]{付録}にまとめた。たとえば命題 $P$ の証明が$p$である時、 この関係は $p:P$（$p$は型$P$の項）という式で表現される。項$p$を構築する厳密な規則が用意されているため、その規則に従っているかを機械的に確認する作業を通して証明の正しさを検証できる。その意味で、コンピュータにとって検証のしやすい形式体系であると言える。

    \end{body}

    \section{証明支援系の概要}
    \begin{body}
        \indent 証明支援系とは、数学の定義・定理・証明を形式言語で記述することに特化したプログラミング言語、かつその言語で記述されたコードが公理や推論規則に照らして正しいかを機械的に検証する機能を備えたもの、の総称である。
        
        \indent 「証明を自動で発見する」ことを目的としたATP（Automated Theorem Proving）と、「与えられた証明が正しいことを厳密に検査する」ことを目的としたITP（Interactive Theorem Proving）が存在する。

        \indent 証明支援系には基礎となっている形式体系の種類に応じて複数の系統がある。代表的なものとして、依存型理論に基づくLeanやCoq（現在Rocqに改名）、高階論理系のIsabelle、HOLなどが存在し、目的に応じて使い分けられている。

        \indent この技術は応用面では、まずソフトウェアの形式検証に用いられる。ソフトウェアはそれぞれ仕様という設計上の基準を持つ。この仕様を数学的な命題として書き下すことで、実装されたプログラムがその仕様を満たしているか否かという問題を、数学的な真偽判定の問題に帰着させることができる。挙動をあらかじめ検証しておくことで、信頼性の高い堅牢なソフトウェア開発が可能となる。

        \indent また近年、証明支援系はLLMとも結びつき始めている。あるLLMがIMO 2025(国際数学オリンピック)において金メダル基準相当の成績に達したとする報告も出ており\cite{DeepMindIMO2025}、自然言語推論の性能が飛躍的に向上している。一方で、もっともらしいが誤りを含む出力（ハルシネーション）は依然としてLLMの課題である。
        しかし、近年注目されている定理証明器を統合した枠組みでは、自然言語の解答案全体または一部を形式化して検証し、失敗時にはフィードバックにより解答を再生成するというループが可能となる。そのため、形式検証が及ぶ範囲についてはハルシネーションを大幅に抑制でき、結果として出力全体に対するハルシネーションの頻度も抑えられる。これは、ハルシネーション自体を抑えるのではなく正確な解答を返すまで出力を繰り返させるという発想の転換である。
        （下図参照、問題文の形式化や自然言語への説明生成には依然として誤りが入りうる）。よって、「形式検証の及ぶ範囲を拡大する」、「問題文の形式化精度を向上させる」ことで生成AIの数学力は向上すると考えられており、研究が盛んに行われている。

    \end{body}

    \begin{tikzpicture}
        %========================
        % 共通
        %========================
        \pgfmathsetlengthmacro{\Span}{0.90*\linewidth}

        \pgfmathsetlengthmacro{\TopY}{0pt}
        \pgfmathsetlengthmacro{\BotY}{-3.5cm}

        \coordinate (TopL) at (-0.5*\Span,\TopY);
        \coordinate (TopR) at ( 0.5*\Span,\TopY);
        \coordinate (BotL) at (-0.5*\Span,\BotY);
        \coordinate (BotR) at ( 0.5*\Span,\BotY);

        %========================
        % 上段
        %========================
        \tikzset{
        box/.style={
        rectangle, draw, thick,
        text width=3.0cm,
        minimum height=1.0cm,
        align=center, font=\small,
        inner sep=2pt
        },
        myarrow/.style={->, >=Stealth, thick, shorten <=2pt, shorten >=2pt},
        danger/.style={box, fill=red!20},
        }

        \node[box] (nlProblem)   at ($(TopL)!0.00!(TopR)$) {自然言語\\問題};
        \node[box] (nlReasoning) at ($(TopL)!0.42!(TopR)$) {自然言語\\推論};
        \node[box] (nlAnswer)    at ($(TopL)!0.84!(TopR)$) {自然言語\\解答};

        \draw[myarrow] (nlProblem) -- (nlReasoning);
        \draw[myarrow] (nlReasoning) -- (nlAnswer);

        \node[above=3mm of nlProblem] {\bfseries 従来の自然言語推論};

        %========================
        % 下段
        %========================
        \tikzset{
        boxS/.style={
        rectangle, draw, thick,
        text width=1.85cm,
        minimum height=0.72cm,
        align=center,
        font=\scriptsize,
        inner sep=2pt
        },
        dangerS/.style={boxS, fill=red!20},
        safeS/.style={boxS, fill=green!20},
        myarrowS/.style={->, >=Stealth, thick, shorten <=2pt, shorten >=2pt}
        }

        \node[boxS]    (flProblemNl) at ($(BotL)!0.00!(BotR)$) {自然言語\\問題};
        \node[safeS]   (flProblem)   at ($(BotL)!0.21!(BotR)$) {形式言語\\問題};
        \node[safeS]   (flReasoning) at ($(BotL)!0.42!(BotR)$) {形式言語\\推論};
        \node[safeS]   (flAnswer)    at ($(BotL)!0.63!(BotR)$) {形式言語\\解答};
        \node[boxS] (finalAnswer) at ($(BotL)!0.84!(BotR)$) {自然言語\\解答};

        \draw[myarrowS] (flProblemNl) -- (flProblem);
        \draw[myarrowS] (flProblem)   -- (flReasoning);
        \draw[myarrowS] (flReasoning) -- (flAnswer);
        \draw[myarrowS] (flAnswer)    -- (finalAnswer);

        \begin{pgfonlayer}{background}
            \node[draw, thick, inner sep=2pt, rounded corners=2pt, fit=(flProblem)(flReasoning)(flAnswer), green, label={[font=\scriptsize]above:ハルシネーションを抑制}] (flBigBox) {};
        \end{pgfonlayer}

        \draw[myarrowS, dashed]
        (flReasoning.south) .. controls +(0,-0.8) and +(0,-0.8) ..
        node[pos=0.55, right, font=\tiny, xshift=-2mm, yshift=-2mm] {検証失敗}
        (flProblem.south);

        \node[above=3mm of flProblemNl] {\bfseries 形式検証統合推論};

    \end{tikzpicture}

    \begin{body}
        \indent 上記を踏まえると本稿の目的は、\emph{数学という体系（特にチェバの定理）を証明支援系を通してコンピュータに検証してもらうために適切な形式言語を選択しコンピュータに伝達すること}である。

        \indent この手段として私はLean4という言語を選択した。
    \end{body}

    \section{Lean4・mathlib4の概要}
    \begin{body}
        \indent Lean4とは、依存型理論の一種CIC（Constructive Inductive Calculus）を基礎理論とする、プログラミング言語であり証明支援系（ITP）でもある。

        \indent この言語の概要を知るにはまず一般的なLean4のファイル構成を把握するのが良いが、これは一般的な数学書と大きく変わらない。論理の流れに従い$def$、$lemma$や$theorem$という宣言に続いて用語の定義、命題の主張が並んでいる。数学書が一部の定理の証明を他の文献に譲るように、Lean4のファイルでは別のファイルを参照し、必要な定理をインポートすることでその定理を既知の定理として証明中で再利用できるようになる。再利用前提でファイルを作成する必要があるため、他のメンバーが理解・共有しやすいようファイル構成にもある程度の形式・ルールが生まれる。この要件を満たしながら定義・定理の数々をLean4に書き起こして行ったファイルを集めたものが、前述したmathlib4という巨大なライブラリである。証明支援系、ライブラリという人によっては馴染みのない用語を聞いて身構える読者もいるかもしれないが、Lean4を一種のプログラミング言語として捉えてしまえば、数学体系をLean4という言語で書き直し、書き直した証明をmathlib4という共用のフォルダにまとめて保存しているだけである。世界中の数学者や技術者が共同でこのフォルダを管理・運営し、世界規模で開発が行われている。この言語は2021年にLean3を改良して生まれた新しい言語であるが、両者に互換性がないため、現在Lean3時代のライブラリ（mathlib3）からmathlib4への翻訳・移管作業が進められている。

        \indent そんなLean4であるが、数学書と最も大きく異なるのがその言語が型理論をベースにして構築されている点である。一般的な数学書は一階述語論理をベースに部分的に自然言語による説明を交えた形式で綴られるが、Lean4ではそのすべてが一つの言語を用いて記述される。数学的主張を世に出す立場として、新たな言語に慣れるまで少し時間がかかるのはデメリットであるが、一方でLean4が持つ大きなメリットが「証明を書く過程でその証明が正しいかどうかを判定してくれる点」である。背後で依存型理論をベースとした型チェッカーというツールが稼働し続け、記述している証明が誤っていれば直ちにエラーを返してくれる場合があり、非常に効率的に記述が行える。いわば、"添削機能付きの数学書"である。

        \indent Lean4における証明記述には大別して2つのモードが存在する。
        
        \begin{bodyitemize}[leftmargin=1.8em]
          \item[] \textbf{項モード}\\
          証明に相当する項をコーディングしながら自分で構成していく。
          \item[] \textbf{tacticモード}\\
          あらかじめ定義された tactic と呼ばれるコマンドを利用して、証明すべき命題を変更するなどの操作を行う。
        \end{bodyitemize}
        

        \vskip\baselineskip

        \indent 詳細な説明は割愛するが、前者は証明全体の流れを明示しやすく，後者は手計算の証明手順に近い形で推論を整理できるため，両者は相補的に用いられる。
    \end{body}

    % =========================================================
    \chapter{構文論における定理の一般化}
    % =========================================================
    \begin{body}
        \indent この章では、一般化とはそもそも何か、また定理の逆とは何かを構文論の観点から振り返る。
    \end{body}

    \section{定理の一般化とは}

    \begin{body}
        \indent 命題の集合$T$に対し、
        \[
        T \vdash \phi
        \]
        の形で表される命題$\phi$のことを（$T$から導出される）\emph{定理}と呼ぶ。特に$\phi$が$p\to q$の形のとき、命題$q\to p$を\emph{定理の逆}と呼ぶ。元の定理が$T$から導出されても、$T$から定理の逆が導出されるとは一般的には言えない。   

        \indent また、複数の命題を論理積でまとめて1つの定理と呼ぶ場合も多い。$\phi_{1},\phi_{2},\cdots,\phi_{n}$が全て$T$から導出される定理である時、
        \[       
        T \vdash (\phi_{1}\land\phi_{2}\land\cdots\land\phi_{n})
        \]
        一方、自然演繹を推論規則として採用する上では演繹定理と呼ばれる以下の性質が成り立つ。
        \[
        T\cup\{\psi\}\vdash \phi \;\Longleftrightarrow\; T\vdash (\psi\to\phi).
        \]
        $T= \{\varphi_{1},\varphi_{2},\cdots,\varphi_{n}\}\cup\{\psi_{1},\psi_{2},\ldots,\psi_{m}\}$の時、
        \[
        \begin{alignedat}{2}
        &\phantom{\iff}\; T \vdash \phi && \\
        &\iff \{\varphi_{1},\varphi_{2},\cdots,\varphi_{n}\}\vdash\bigl(\psi_{1}\to(\psi_{2}\to(\cdots\to(\psi_{m}\to\phi)))\bigr) &&\text{\footnotesize（演繹定理を繰り返し適用）}\\
        &\iff \{\varphi_{1},\varphi_{2},\cdots,\varphi_{n}\}\vdash\bigl(\neg\psi_{1}\lor(\neg\psi_{2}\lor(\cdots\lor(\neg\psi_{m}\lor\phi)))\bigr) &&\text{\footnotesize（含意と論理和の関係）}\\
        &\iff \{\varphi_{1},\varphi_{2},\cdots,\varphi_{n}\}\vdash\bigl(\neg\psi_{1}\lor\neg\psi_{2}\lor\cdots\lor\neg\psi_{m}\lor\phi\bigr) &&\text{\footnotesize（$\lor$の結合法則）}\\
        &\iff \{\varphi_{1},\varphi_{2},\cdots,\varphi_{n}\}\vdash\bigl(\neg(\psi_{1}\land\psi_{2}\land\cdots\land\psi_{m})\lor\phi\bigr) &&\text{\footnotesize（ド・モルガンの法則）}\\
        &\iff \{\varphi_{1},\varphi_{2},\cdots,\varphi_{n}\}\vdash\bigl((\psi_{1}\land\psi_{2}\land\cdots\land\psi_{m})\to\phi\bigr) &&\text{\footnotesize（含意と論理和の関係）}\\
        \end{alignedat}
        \]
        よって、$\phi_{1},\phi_{2},\cdots,\phi_{n}$が$T$から導出される定理である時、以下のように書き下せる。
        \[
        \{\varphi_{1},\varphi_{2},\cdots,\varphi_{n}\}\vdash{\psi_{1}\land\psi_{2}\land\cdots\land\psi_{m}}\to{\phi_{1}\land\phi_{2}\land\cdots\land\phi_{n}}
        \]
        
        \vskip\baselineskip

        $P:=\{\varphi_{1},\varphi_{2},\cdots,\varphi_{n}\},A:=\{\psi_{1},\psi_{2},\cdots,\psi_{m}\},C:=\{\phi_{1},\phi_{2},\cdots,\phi_{n}\}$と置き、便宜上以下のように略記する
        \[
        P\vdash A\to C
        \]
        これは自然演繹の公式の記法ではなく、説明上の都合でこのような書き方になっていることに留意してほしい。
        この$P$を\emph{前提/背景理論}、$A$を\emph{仮定}、$C$を\emph{結論}と呼ぶ。
        要するに、定理が与えられるとそれを適当な粒度の命題の組を用意し、\emph{前提・仮定・結論の導出関係として書き換えることができる}ということである。

        \indent 数学基礎論以外の分野の一般的な数学書では前提$P$として$ZF/ZFC$が慣習的に用いられるため、これらは暗黙に受け入れたものとして敢えては明記されず、仮定$A$から結論$C$を導く、というステップのみが記述されることが多い。

        \indent また、この形式で定理を表現することにより、暗黙に認めていた前提を明文化できるため、Lean4での形式化に繋げやすいというメリットがある。

        \indent この上で、本稿で言う「定理の一般化」とは、既存の定理がある意味で\emph{特別な場合として含まれる}ような、より広い主張をする定理を与えることである。典型的には次の三方向がある。
        \begin{bodyitemize}[leftmargin=3.5em]
            \item \textbf{前提の弱化}：$P$の元の命題の主張を弱める。$\varphi\in P$を、$\varphi\to \varphi'$を満たす$\varphi'$で置き換える。 
            \item \textbf{仮定の弱化}：$A$の元の命題の主張を弱める。$\psi\in A$を、$\psi\to \psi'$を満たす$\psi'$で置き換える。
            \item \textbf{結論の強化}：$C$の元の命題の主張を強める。$\phi\in C$を、$\phi'\to \phi$を満たす$\phi'$で置き換える。
            \item \textbf{複数の組合せ}：上記を複数個組み合わせる。
        \end{bodyitemize}
        
        \indent 3つ目は「定理の強化」と呼ばれることが多いが、本稿では一般化に含める。一般化と聞くと仮定と結論に着目しがちであるが、定理が依存する前提を緩め（例：体から環へ、可換環から一般の環へ、有限次元から任意次元へ）、対象の主張をより広い土台で成立させることも一般化の重要な形である。

    \end{body}

    % =========================================================
    \chapter{チェバの定理とアフィン空間}
    % =========================================================
    \begin{body}
        \indent チェバの定理とは高校数学でも頻繁に登場する、初等幾何における重要な定理の一つである。

        \indent 本章ではこのチェバの定理の詳細と、チェバの定理が展開されるアフィン空間の性質について議論する。

    \end{body}

    \section{チェバの定理の有名な形}

    \begin{body}
        \indent 高校の教育課程において、チェバの定理は以下の形で紹介されることが多い。ただし、定理中で登場する用語は初等幾何的な高校の学習課程内の意味である。
    \end{body}

    \begin{theorem}[古典的チェバの定理]\label{thm:classical-ceva}
        三角形$\triangle ABC$の各辺$BC,CA,AB$またはその延長線上にそれぞれ頂点とは異なる点$D,E,F$をとる。
        このとき$3$直線$AD$、$BE$、$CF$が$1$点で交わるならば、以下の式が成立する
        \[
          \dfrac{BD}{DC}\,\dfrac{CE}{EA}\,\dfrac{AF}{FB} = 1
        \]
    \end{theorem}

    \begin{theorem}[古典的チェバの定理の逆]
        三角形$\triangle ABC$の各辺$BC,CA,AB$またはその延長線上にそれぞれ頂点とは異なる点$D,E,F$をとる。
        このとき
        \[
          \dfrac{BD}{DC}\,\dfrac{CE}{EA}\,\dfrac{AF}{FB} = 1
        \]
        が成立するならば、$3$直線$AD,BE,CF$は$1$点で交わる、または互いに平行である。
    \end{theorem}

    \begin{body}
        \indent ここで登場する点$D,E,F$を\emph{チェビアンの足}、$3$直線$AD,BE,CF$を\emph{チェビアン}と呼ぶ。

        \indent つまり古典的チェバの定理（またはその逆）とは\emph{チェビアンが共点するという幾何的性質が成立する条件を代数的に記述する定理}である。

        \indent 上記のチェバの定理の主張には、一見すると線分の長さ（BD、DC等）が現れるため、ユークリッド幾何に依存する定理のように見える。

        \indent しかし、実際にはこの定理の主張は線分の長さの比（$\dfrac{BD}{DC}$ 等）のみで記述されている。実はこの概念は後述するようにユークリッド幾何を用いることなくアフィン幾何という体系内で定義することができる。

        \indent またチェバの定理が必要十分性を主張する定理だと認識していた読者も多いかもしれない。$D,E,F$を辺上に限定する版や、$3$直線が$1$点で交わる必要十分条件を提供する版など、教科書により紹介の仕方に微妙な違いがあり、これらは教育的配慮によるものである。詳細は\hyperref[appendix:classical-ceva-version]{付録}で解説するが、アフィン幾何の範疇では古典的チェバの定理の主張が「共点$⇔$等式条件」というシンプルな形の必要十分条件では表せないためこのような紹介になっているという事にだけ留意してほしい。
    \end{body}


    \section{概念の定義}
    \begin{body}
        \indent まず今後の議論で用いる概念を定義する。ここでは紙面上の証明に直接使いやすい形の定義を採用している。mathlib4 で採用されている実装との対応は後の章で説明する。

        \indent 以下では $k$ を可換環、$V$ を $k$-加群、$P$ を $V$ 上のアフィン空間、$\iota$ を有限添字集合とする。
    \end{body}

    \begin{definition}[アフィン空間]
        $P$ が $V$ 上の（$k$ に関する）\emph{アフィン空間}であるとは、次の2つの写像が存在し、以下の公理を満たすことをいう：
        \[
        \dotplus : V\times P \ni (v,p) \mapsto v \dotplus p \in P,
        \qquad
        \dotminus : P\times P \ni (p,q) \mapsto p \dotminus q \in V.
        \]
        \begin{theoremitemize}[leftmargin=5em, topsep=2pt, partopsep=0pt, itemsep=0pt, parsep=0pt]
            \item 任意の $p\in P$ と $v,w\in V$ について
            \[
              \mathbf 0 \dotplus p = p,
              \qquad
              (v+w)\dotplus p = v\dotplus (w\dotplus p).
            \]
            \item 任意の $p,q\in P$ について
            \[
              (p\dotminus q)\dotplus q = p.
            \]
            \item 任意の $p\in P$ と $v\in V$ について
            \[
              (v\dotplus p)\dotminus p = v.
            \]
        \end{theoremitemize}
        アフィン空間 $P$ の元を\emph{点}と呼ぶ。写像 $p:\iota\to P$ またはその像を\emph{点族}と呼ぶ。
    \end{definition}

    \begin{remark}[点に対する演算について]
        アフィン空間で最初から定義されているのは「点同士の差からベクトルを作る」演算と「ベクトルで点を平行移動する」演算であり、点どうしの和 $p+q$ や点のスカラー倍 $\lambda p$ は一般には定義されない。
        したがって以降の「点の線形結合」に見える式は、厳密にはアフィン結合の略記として理解する必要がある。
    \end{remark}

    \begin{abbreviation}[記法]
        以下、簡単のため $v\dotplus p$ を $v+p$、$p\dotminus q$ を $p-q$ と書く。
    \end{abbreviation}

    \begin{definition}[アフィン結合（\texorpdfstring{$\AffComb$}{AffComb}）]
        $s\subseteq \iota$ を空でない有限集合とし、重み $\lambda:s\to k$ が
        \[
          \sum_{i\in s}\lambda(i)=1
        \]
        を満たすとする。このとき任意に $j\in s$ を固定して
        \[
          \AffComb\bigl(p|_s,\lambda\bigr)
          :=
          \left(\sum_{i\in s}\lambda(i)\,(p(i)-p(j))\right)+p(j)
        \]
        と定める。この表現を（点族 $p|_s$ の）\emph{アフィン結合}と呼び、$\lambda$ を\emph{重み}と呼ぶ。
    \end{definition}

    \begin{remark}[アフィン結合の基本性質]
        条件 $\sum_{i\in s}\lambda(i)=1$ のもとでは、上の定義は基点 $j\in s$ の選び方に依らない。したがって形式的な略記として
        \[
          \AffComb\bigl(p|_s,\lambda\bigr)=\sum_{i\in s}\lambda(i)\,p(i)
        \]
        と書いてよい。

        また、2点 $A,B\in P$ に対して $q(0)=A,\ q(1)=B$ とおくと、$\{0,1\}$ 上のアフィン結合は
        \[
          \AffComb(q,(1-\tau,\tau)) = \tau(B-A)+A
          \qquad (\tau\in k)
        \]
        と書ける。したがって 2 点を結ぶ直線は 1 つのスカラー・パラメータで記述できる。
    \end{remark}

    \begin{definition}[アフィン包]
        \[
          \AffHull(p,s)
          :=
          \left\{\,
            \AffComb\bigl(p|_s,\lambda\bigr)
            \ \middle|\
            \lambda:s\to k,\ \sum_{i\in s}\lambda(i)=1
          \,\right\}
        \]
        で定まる集合を点族 $p|_s$ の \emph{アフィン包}と呼ぶ。
    \end{definition}

    \begin{definition}[アフィン部分空間]
        部分集合 $S\subseteq P$ が \emph{アフィン部分空間}であるとは、$S\neq\emptyset$ であり、任意の有限部分集合 $t\subseteq S$ と任意の重み $\mu:t\to k$ について
        \[
          \sum_{x\in t}\mu(x)=1
          \quad\Longrightarrow\quad
          \AffComb(t,\mu)\in S
        \]
        が成り立つことをいう。
    \end{definition}

    \begin{remark}[最小性]
        $\AffHull(p,s)$ は $p(s)$ を含む最小のアフィン部分空間である。すなわち、$p(s)\subseteq T\subseteq P$ を満たす任意のアフィン部分空間 $T$ に対し
        \[
          \AffHull(p,s)\subseteq T
        \]
        が成り立つ。
    \end{remark}

    \begin{definition}[直線]
        異なる 2 点 $A,B\in P$ に対し、$q:\{0,1\}\to P$ を $q(0)=A,\ q(1)=B$ として
        \[
          \mathrm{Line}(A,B):=\AffHull(q,\{0,1\})
        \]
        とおく。これを $A$ と $B$ を通る\emph{直線}と呼ぶ。
    \end{definition}

    \begin{remark}[直線上の点の表示]
        上の記法のもとで
        \[
          \mathrm{Line}(A,B)=\{\,\tau(B-A)+A\mid \tau\in k\,\}
        \]
        が成り立つ。よって「点 $X$ が $A,B$ を通る直線上にある」とは、ある $\tau\in k$ が存在して $X=\tau(B-A)+A$ と書けることと同値である。
    \end{remark}

    \begin{definition}[平行]
        アフィン部分空間 $S,T\subseteq P$ に対し、$T$ が $S$ に \emph{平行}であるとは、ある $v\in V$ が存在して
        \[
          T=\{\,v+p\mid p\in S\,\}
        \]
        と書けることをいう。このとき $S\parallel T$ と書く。
    \end{definition}

    \begin{definition}[アフィン独立]
        空でない有限集合 $s\subseteq \iota$ に対し、点族 $p|_s$ が \emph{アフィン独立}であるとは、ある（したがって任意の） $i_0\in s$ に対して、族
        \[
          \{\,p(i)-p(i_0)\mid i\in s\setminus\{i_0\}\,\}
        \]
        が $V$ で線形独立であることをいう。
    \end{definition}

    \begin{remark}[アフィン独立から従う事実]
        $p|_s$ がアフィン独立であるとき、以下が成り立つ。
        \begin{theoremitemize}[leftmargin=5em, topsep=2pt, partopsep=0pt, itemsep=0pt, parsep=0pt]
            \item 任意の $q\in \AffHull(p,s)$ は、ただ一つの重み $\lambda:s\to k$ を用いて
            \[
              q=\AffComb\bigl(p|_s,\lambda\bigr),
              \qquad
              \sum_{i\in s}\lambda(i)=1
            \]
            と書ける。すなわち重心座標は一意である。
            \item 任意の $i_0\in s$ と係数族 $\lambda:s\to k$ について
            \[
              \sum_{i\in s}\lambda(i)=0,
              \qquad
              \sum_{i\in s}\lambda(i)\,(p(i)-p(i_0))=0
            \]
            ならば、全ての $i\in s$ で $\lambda(i)=0$ である。特に、総和が $0$ である係数族が自明でなければ、対応するベクトル和は $0$ にならない。
        \end{theoremitemize}
    \end{remark}

    \begin{definition}[重心座標]
        $p|_s$ がアフィン独立であるとする。任意の点 $q\in \AffHull(p,s)$ に対し、ただ一つの重み $\lambda:s\to k$ が存在して
        \[
          q=\AffComb\bigl(p|_s,\lambda\bigr),
          \qquad
          \sum_{i\in s}\lambda(i)=1
        \]
        を満たす。この $\lambda$ を $q$ の（$p|_s$ に関する）\emph{重心座標}と呼ぶ。
    \end{definition}

    \begin{definition}[単体]
        アフィン独立な点族 $p|_s$ を \emph{$(|s|-1)$-単体}と呼ぶ。
    \end{definition}

    \begin{example}
        $s$ が $n+1$ 点集合のとき、アフィン独立な $p|_s$ を $n$-単体と呼ぶ。特に異なる 3 点 $A,B,C$ がアフィン独立なら、それは $2$-単体であり、通常の三角形 $\triangle ABC$ に対応する。
    \end{example}

    \begin{remark}
        教科書では単体を「内部や境界まで含めた図形」として述べることが多いが、mathlib4 ではまず「アフィン独立な頂点族」として単体を表現し、その上で対応する凸包や内部を別に扱う。
    \end{remark}

    \begin{definition}[面]
        $S$ を $p|_s$ が定める単体とする。$s_0\subseteq s$ に対し、点族 $p|_{s_0}$ が定める単体を $S$ の \emph{面}と呼ぶ。特に $|s_0|=k+1$ のとき、その面を \emph{$k$-面}と呼ぶ。
    \end{definition}

    \begin{definition}[対向面]
        $S$ を $p|_s$ が定める単体とし、$i\in s$ とする。$s\setminus\{i\}$ が定める面を、頂点 $p(i)$ に対する \emph{対向面}と呼ぶ。
    \end{definition}

    \begin{definition}[チェビアン]
        $S$ を $p|_s$ が定める単体とし、$i\in s$ とする。対向面のアフィン包 $\AffHull(p,s\setminus\{i\})$ 上の点 $q$ をとる。このとき、頂点 $p(i)$ と $q$ を結ぶ直線 $\mathrm{Line}(p(i),q)$ を、頂点 $p(i)$ からの \emph{チェビアン}と呼ぶ。点 $q$ はそのチェビアンの\emph{足}に相当する。
    \end{definition}

    \begin{remark}[対向面と重心座標]
        $p|_s$ がアフィン独立で、$q\in \AffHull(p,s)$ の重心座標を $\lambda:s\to k$ とする。すると
        \[
          q\in \AffHull(p,s\setminus\{i\})
          \quad\Longleftrightarrow\quad
          \lambda(i)=0
        \]
        が成り立つ。したがって「対向面上にある」という条件は、対応する頂点の重心座標が消えることとして表現できる。
    \end{remark}

    \begin{definition}[共点する]
        直線 $S_1,\dots,S_m$ が 1 点 $x\in P$ を共有するとき、すなわち
        \[
          x\in \bigcap_{\ell=1}^m S_\ell
        \]
        を満たすとき、これらは \emph{共点する}という。また $x$ を\emph{共点}と呼ぶ。
    \end{definition}

    \section{古典的チェバの定理の書き換え}

    \begin{body}
        \indent 上記の概念を用いて古典的チェバの定理及びその逆を導出関係を明示しながら書き換える。一般化を見据えて冗長な表現になっていることに留意してほしい。
    \end{body}

    \begin{theorem}[古典的チェバの定理の言い換え]

        \noindent\textbf{前提：}\par
        \begin{theoremenumerate}[label=(P\arabic*), leftmargin=4em, topsep=2pt, partopsep=0pt, itemsep=0pt, parsep=0pt]
            \item $k$ は可換体である。
            \item $V$ は $k$-ベクトル空間である。
            \item $P$ は $V$ 上のアフィン空間である。
            \item $\iota = \{0,1,2\}$ である。
            \item 点族$p:\iota \to P$ はアフィン独立で、$A:=p(0), B:=p(1), C:=p(2)$と表す。
            \item $s = \iota$          
            \item 各 $i\in s$ について、有限集合 $F_i:=\iota\setminus\{i\}\subset \iota$ と重み $w_i:F_i\to k$が与えられ、
            \[
            \sum_{j\in F_i} w_i(j)=1
            \]
            が成立し、さらに $t,u,v \in k \setminus \{0,1\}$ であって、$D:=(1-t)B+tC, E:=(1-u)C+uA, F:=(1-v)A+vB$ と表される。
        \end{theoremenumerate}
        
        \noindent\textbf{仮定：}\par
        \begin{theoremenumerate}[label=(A\arabic*), leftmargin=4em, topsep=2pt, partopsep=0pt, itemsep=0pt, parsep=0pt]
            \item ある点 $p' \in P$ が存在し、$p'$は直線$AD, BE, CF$上に存在する。
        \end{theoremenumerate}
        
        \noindent\textbf{結論：}\par
        \begin{theoremenumerate}[label=(\roman*), leftmargin=4em, topsep=2pt, partopsep=0pt, itemsep=0pt, parsep=0pt]
            \item 以下の等式が成立する。
            \[
            \left(\frac{t}{1-t}\right)\left(\frac{u}{1-u}\right)\left(\frac{v}{1-v}\right)=1
            \]
        \end{theoremenumerate}
    \end{theorem}

    \begin{body}
        \indent この言い換えが適切に古典的チェバの定理の言い換えになっていることを確かめる。

        \indent $k$が可換体なのは等式条件に商が含まれていること、積の順序を変えても成り立つであろうことから明らかである。$P$が$k$-ベクトル空間上のアフィン空間であることは後に示す。$\iota = \{0,1,2\}$ であることは三角形（$3$点からなる点族）を考えていることを意味し、アフィン独立性は$3$頂点が同一直線上となり三角形が潰れてしまうことを防ぐためである。$2$頂点を結ぶ直線の重心座標は一つのパラメータを用いて表せる。
    \end{body}
    
    \begin{theorem}[古典的チェバの定理の逆の言い換え]

        \noindent\textbf{前提：}\par
        \begin{theoremenumerate}[label=(P\arabic*), leftmargin=4em, topsep=2pt, partopsep=0pt, itemsep=0pt, parsep=0pt]
            \item $k$ は可換体である。
            \item $V$ は $k$-ベクトル空間である。
            \item $P$ は $V$ 上のアフィン空間である。
            \item $\iota = \{0,1,2\}$ である。
            \item 点族$p:\iota \to P$ はアフィン独立で、$A:=p(0), B:=p(1), C:=p(2)$と表す。
            \item $s = \iota$          
            \item 各 $i\in s$ について、有限集合 $F_i:=\iota\setminus\{i\}\subset \iota$ と重み $w_i:F_i\to k$が与えられ、
            \[
            \sum_{j\in F_i} w_i(j)=1
            \]
            が成立し、さらに $t,u,v \in k \setminus \{0,1\}$ であって、$D:=(1-t)B+tC, E:=(1-u)C+uA, F:=(1-v)A+vB$ と表される。
        \end{theoremenumerate}
    
        \noindent\textbf{仮定：}\par
        \begin{theoremenumerate}[label=(A\arabic*), leftmargin=4em, topsep=2pt, partopsep=0pt, itemsep=0pt, parsep=0pt]
            \item 係数 $t,u,v$ が
            \[
            \left(\frac{t}{1-t}\right)\left(\frac{u}{1-u}\right)\left(\frac{v}{1-v}\right)=1
            \]
            を満たす。
        \end{theoremenumerate}
    
        \noindent\textbf{結論：}\par
        \begin{theoremenumerate}[label=(\roman*), leftmargin=4em, topsep=2pt, partopsep=0pt, itemsep=0pt, parsep=0pt]
            \item 点$p' \in P$が存在して$p'$ は直線 $AD, BE, CF$ 上に存在する、または直線 $AD, BE, CF$ が互いに平行である。
        \end{theoremenumerate}
    \end{theorem}

    \begin{body}
        \indent 前提条件は順方向と同じである。
    \end{body}

    \section{チェバの定理とアフィン幾何学}
        
    \begin{body}
        \indent 先ほど紹介したように、この定理の舞台はアフィン空間であり、アフィン幾何学の定理として数えられる。ある定理がアフィン幾何学の定理として主張するには「その定理がアフィン変換で不変な量だけを使って定式化でき、しかも命題の真偽がアフィン変換で保存される」ことを示せば良い。

        \indent 上記の定理は全てアフィン空間の概念で記述されており、線分の長さをはじめとしたユークリッド幾何の概念は一切登場しない。辺の長さの比に見えていたものは重心座標で現れる係数の比であり、
        \[
        \dfrac{BD}{DC} := \dfrac{t}{1-t},\qquad
        \dfrac{CE}{EA} := \dfrac{u}{1-u},\qquad
        \dfrac{AF}{FB} := \dfrac{v}{1-v}
        \]
        である。

        \indent この結果は本稿で作成するモジュールをmathlib4のどのディレクトリのファイルに追加するか決定する際に重要となる。一般化を行うときもユークリッド幾何の計量の概念は避けることが必要である。

    \end{body}

    \begin{theorem}[古典的チェバの定理のアフィン性]
        前節で与えた古典的チェバの定理およびその逆の言い換えは、可逆なアフィン変換で保存される概念のみを用いて定式化されている。したがって、それらはアフィン幾何の定理である。
    \end{theorem}

    \begin{proof}
        前節の定式化に現れるのは、アフィン独立性、アフィン結合、アフィン包、直線への所属、そして係数 $t,u,v$ の間の代数的等式のみである。これらはいずれも長さ・角度・内積といった計量概念を用いない。

        実際、可逆なアフィン変換 $f:P\to P$ はアフィン結合を保存するので
        \[
          D=(1-t)B+tC
          \quad\Longrightarrow\quad
          f(D)=(1-t)f(B)+t f(C)
        \]
        が成り立ち、$E,F$ についても同様である。したがって、$t,u,v$ はユークリッド距離からではなく重心座標として定まる量であり、アフィン変換で不変である。

        また、点のアフィン包への所属はアフィン変換で保存されるから、直線 $AD,BE,CF$ の共点性も保存される。逆方向で現れる「互いに平行である」という条件も、直線の方向部分空間が一致することとして述べられるので、やはりアフィン変換で保存される。

        以上より、古典的チェバの定理およびその逆は、いずれもアフィン幾何の範疇で自然に述べられる定理である。
    \end{proof}

    % =========================================================
    \chapter{チェバの定理の一般化と選定}
    % =========================================================
    \begin{body}
        \indent 本章では上記の議論を踏まえ、一般化チェバの定理の主張として求められる何が相応しいかを明確にする。
    \end{body}


    \section{一般化の方針}

    \begin{body}
        \indent 古典的チェバに対して上記の定式化を行うと、定理がさまざまな方向に一般化できることが見える。

        \indent 代表的には次のような観点がある。
        \begin{bodyitemize}
            \item 周辺空間だけを一般化する。
            \item 三角形を $n$-単体へ一般化する。
            \item 比の積等式を、重心座標や行列式によるより対称的な条件へ置き換える。
            \item 各頂点から 1 本ずつではなく、より多くのチェビアン族を同時に扱う。
        \end{bodyitemize}

        \vskip\baselineskip

        \indent 実際に候補となる主張を並べると次のようになる。

        \begin{bodydescription}
            \item[\hspace*{1.5em}・候補A：周囲空間のみ一般化]
            三角形そのものは保ったまま、周囲空間を平面に限らず任意次元のアフィン空間へ広げる。主張の本質は古典的チェバと同じであり、既存の三角形版を高次元空間に埋め込んで扱う立場である。

            \item[\hspace*{1.5em}・候補B：$n$-単体版]
            $n$-単体の各頂点 $p(i)$ に対し、その対向面のアフィン包上の点 $q(i)$ を取り、直線 $p(i)q(i)$ を考える。これらが共点または平行になる条件を、各 $q(i)$ の重心座標の整合性として述べる定理である。古典的チェバを最も自然に高次元化した版と言える。

            \item[\hspace*{1.5em}・候補C：行列式・体積比による版]
            三角形での「線分比の積」を、高次元では行列式・符号付き体積・質量分布条件のような量に置き換えて述べる版である。幾何的直観は強い一方、mathlib4 上では補助理論がやや重くなる。

            \item[\hspace*{1.5em}・候補D：multipede 型の版]\cite{multipede}
            1 つの頂点から複数の対向面へ向かうチェビアンを同時に考え、チェビアンの本数を頂点数より多く許す版である。一般性は高いが、そのぶん定式化と補助概念が大きくなる。
        \end{bodydescription}

        \indent このような、さまざまな版で一般化されたチェバの定理の総称を \emph{チェバ型定理} と呼ぶ。
    \end{body}


    \section{先行プロジェクトの紹介}
    \begin{body}
        \indent 過去のいくつかのチェバ型定理形式化に関する先行研究を紹介する。

        \begin{bodydescription}
            \item[\hspace*{1.5em}・Lean3 時代の形式化]
            期間: $2021$年$12$月 〜 $2023$年$7$月\\
            GitHub: \url{https://github.com/leanprover-community/mathlib3/pull/10632} \cite{mathlib3_ceva}

            Mantas Bak\v{s}ys 氏による試みである。三角形版のチェバの定理を主に距離を用いて形式化したが、後の議論では「チェバは本質的にはアフィン幾何の定理として実装すべきである」という方向性が共有され、mathlib4 には移植されなかった。

            \item[\hspace*{1.5em}・mathlib4 における一般化版の導入]
            期間: $2026$年$1$月\\
            GitHub: \url{https://github.com/leanprover-community/mathlib4/pull/33409} \cite{mathlib4_ceva_myers}

            Joseph Myers 氏による試みである。アフィン空間上の $n$-単体に対する一般化チェバの定理の順方向を形式化し、三角形版の古典的定理をその系として導いた。2026 年 1 月 14 日に mathlib4 へマージされ、本稿執筆時点で公開ライブラリに入っているチェバ型定理の中心的結果である。
        \end{bodydescription}

        \indent 本稿の位置づけは、この既存の順方向定理を出発点にして、その逆方向を「共点または平行」というアフィン幾何に自然な形で定式化し、形式化に必要な補題を整理することにある。
    \end{body}

    \section{本稿で扱う一般化の選定}
    \begin{body}
        \indent 本稿では既存のmathlib4に用意された一般化チェバの定理の順方向の方針を採用する。

        \indent 上述したJoseph Myers氏によって形式化されたものであり、「各頂点と、その対向面のアフィン包上の点を結ぶ直線群が共点するときに成り立つ必要条件」を与える。
    \end{body}


    \subsection{先行研究の主張}

    \begin{body}
        \indent 現在の mathlib4 には \nolinkurl{Mathlib/LinearAlgebra/AffineSpace/Ceva.lean} が存在し、主な宣言として次の 4 つが公開されている\cite{mathlib4_ceva_docs}。
    \end{body}

    \begin{bodydescription}
        \item[\hspace*{1.5em}・\nolinkurl{AffineIndependent.exists_affineCombination_eq_smul_eq}]
        任意添字集合版の一般化チェバ定理。各頂点と、そこから見たアフィン結合の点を結ぶ直線が共点するとき、共点自身もまた全頂点のアフィン結合で書け、その重みが各局所的な重みと第 $i$ 成分を除いて比例することを述べる。

        \item[\hspace*{1.5em}・\nolinkurl{AffineIndependent.exists_affineCombination_eq_smul_eq_of_fintype}]
        上の結果の有限添字集合版であり、本稿で扱う逆方向定理に最も近い既存結果である。

        \item[\hspace*{1.5em}・\nolinkurl{Affine.Triangle.prod_eq_prod_one_sub_of_mem_line_point_lineMap}]
        三角形版の古典的チェバで、係数の積の等式
        \[
          \prod_i r_i = \prod_i (1-r_i)
        \]
        を与える。

        \item[\hspace*{1.5em}・\nolinkurl{Affine.Triangle.prod_div_one_sub_eq_one_of_mem_line_point_lineMap}]
        さらに体上では、上の等式を
        \[
          \prod_i \frac{r_i}{1-r_i}=1
        \]
        という教科書的な形に書き換えた版も用意されている。
    \end{bodydescription}

    \begin{body}
        \indent すなわち、最初の 2 つが一般化された順方向チェバ型定理であり、後の 2 つはそこから導かれる三角形版の系である。

        \indent 本稿のテーマに最も近いのは 2 番目の有限添字版である。その内容を自然言語で要約すると、次のようになる。
    \end{body}

    \begin{theorem}[チェバ型定理順方向]
        $n$-単体の頂点を$p(0)$,…,$p(n)$とする。
        いくつかの頂点$p(i)$に対し、$p(i)$を含まない任意の面上に点$q(i)$を固定し、直線$p(i)q(i)$を考える。
        こうして得られる複数の直線がある点$p'$で交わる時、$p'$は$p(F_i)\cup \{p(i)\}$が張るアフィン部分空間中に存在して、各iについて$q(i)$の$F_i$が定める点族に関する重心座標、および$p'$の$\iota$が定める点族に関する重心座標は第$i$成分を除いて比例する。
    \end{theorem}

    \begin{body}
        \indent もはや古典的チェバの定理の原形を留めていない。

        \indent 点$q(i)$、直線$p(i)q(i)$それぞれは古典的チェバの定理においてチェビアンの足だったもの、チェビアンだったものである。要するに、「共点が存在すれば、共点とチェビアンの足の重心座標成分の大部分が比例する」ことを主張している。

        \indent 一見古典的チェバとの関連は見えづらい。しかし、以下の点で一般化になっている。

        \indent 一つ目に、\emph{$n$-単体についての主張になっている}。前提の弱化である。$n=2$とすれば三角形についての定理になる。

        \indent 二つ目に、\emph{チェビアンは全ての頂点から下ろす必要がない}。前提の弱化である。直線$p(i)q(i)$がチェビアンに相当する。古典的チェバの定理では、全ての頂点からチェビアンを下ろしている。

        \indent 三つ目に、\emph{チェビアンの足は対向面上でなくて良い}。前提の弱化である。点$q(i)$がチェビアンの足に相当する。古典的チェバの定理では、チェビアンの足は対向面（対辺）上にあった。しかし、この定理中ではチェビアンの足はいずれかの面上にあれば良い。

        \indent 四つ目に、\emph{商の等式条件が係数の比例条件になっている}。結論の強化である。これによりスカラーを体の元に限る必要がなくなる。

        \indent この主張を導出関係を明確にして書き下す。

    \end{body}

    \begin{theorem}[チェバ型定理順方向の言い換え]
        \noindent\textbf{前提：}\par
        \begin{theoremenumerate}[label=(P\arabic*), leftmargin=4em, topsep=2pt, partopsep=0pt, itemsep=0pt, parsep=0pt]
            \item $k$ は 環である。
            \item $V$ は $k$-加群である。
            \item $P$ は $V$ 上のアフィン空間である。
            \item $\iota$ は有限の添字集合である。
            \item 点族 $p : \iota \to P$はアフィン独立。
            \item $s \subseteq \iota$ は空でない。
            \item 各 $i\in s$ について，有限集合 $F_i\subseteq \iota$ と重み $w_i:F_i\to k$ が与えられ，
            \[
            i\notin F_i \land \sum_{j\in F_i} w_i(j)=1
            \]
            が成り立つ。
        \end{theoremenumerate}
        \noindent\textbf{仮定：}\par
        \begin{theoremenumerate}[label=(A\arabic*), leftmargin=4em, topsep=2pt, partopsep=0pt, itemsep=0pt, parsep=0pt]
            \item ある点 $p'\in P$ が存在し、各 $i \in s$ について、$p'$ は2点 $p(i)$ と $\sum_{j \in F_i} w_i(j)p(j)$ を結ぶ直線上に存在する。
        \end{theoremenumerate}

        \noindent\textbf{結論：}\par
        \begin{theoremenumerate}[label=(C\arabic*), leftmargin=4em, topsep=2pt, partopsep=0pt, itemsep=0pt, parsep=0pt]
            \item ある重み $w' : \iota \to k$が存在して
            \[
            \sum_{j \in \iota} w'(j) = 1
            \]
            を満たし、
            \[
            p' = \sum_{j \in \iota} w'(j)p(j)
            \]
            と書ける。
            \item 各 $i \in s$ ごとにスカラー $r_i \in k$ が存在し、次が成り立つ： 
            \[
            \forall i \in s,\ \exists r_i \in k,\ \forall j \in F_i,\ \bigl(w'(j) = r_i\, w_i(j)\bigr).
            \]
        \end{theoremenumerate}
    \end{theorem}

    \begin{proof}  
        (A1) より、任意の $i\in s$ に対して $p'$ は2点 $p(i)$ と $\sum_{j \in F_i} w_i(j)p(j)$ を結ぶ直線上に存在する。
        よって、
        \[
        \exists \alpha_i\in k, p'=\alpha_i\,p(i)+(1-\alpha_i)\sum_{j\in F_i} w_i(j)\,p(j)
        \]
        ここで各 $i\in s$ について、関数 $\bar w_i:\iota\to k$ を次で定める：
        \[
        \bar w_i(j):=
        \begin{cases}
        \alpha_i & (j=i),\\
        (1-\alpha_i)\,w_i(j) & (j\in F_i),\\
        0 & (j\in\iota\setminus(F_i\cup\{i\})).
        \end{cases}
        \]
        （$w_i: F_i\to k$ であり、$\iota$全域で定義されていないため、$j\notin F_i$ の場合を上の定義により$\bar w_i(j)=0$ として埋めた。）
  
        $i\notin F_i$ と $\sum_{j\in F_i}w_i(j)=1$ より
        \[
        \begin{aligned}
        \sum_{j\in\iota} \bar w_i(j)
            &= \bar w_i(i)+\sum_{j\in F_i} \bar w_i(j)+\sum_{j\in\iota\setminus(F_i\cup\{i\})} \bar w_i(j)\\
            &= \alpha_i + (1-\alpha_i)\sum_{j\in F_i} w_i(j) + 0\\
            &= \alpha_i + (1-\alpha_i)\\
            &= 1
        \end{aligned}
        \]
        が成り立つ。また上の定義から
        \[
        p'=\sum_{j\in\iota} \bar w_i(j)\,p(j)
        \]
  
        以上より、各 $i\in s$ に対して
        \[
        \sum_{j\in\iota} \bar w_i(j)=1,\qquad
        p'=\sum_{j\in\iota} \bar w_i(j)\,p(j)
        \]
        が得られた。
  
        また(P5)より、$p:\iota\to P$はアフィン独立であるから$\AffHull(p)$中の点を点族$p$のアフィン結合で一意に表せる。
        (A1)より、$p'\in P$である。
        従って任意の $i,i'\in s$ について
        \[
        \bar w_i = \bar w_{i'}
        \]
  
        よって、任意に $i_0\in s$ を一つ固定して
        \[
        w' := \bar w_{i_0}:\iota\to k
        \]
        とおけば、全ての $i\in s$ に対して $w'=\bar w_i$。
  
        したがって
        \[
        \sum_{j\in\iota} w'(j)=1,\qquad
        p'=\sum_{j\in\iota} w'(j)\,p(j)
        \]
        (C1) が示された。
  
        また、$r_i:=1-\alpha_i$ とおけば、定義より$w'=\bar w_i$ なので、任意の $j\in F_i$ について
        \[
        \begin{aligned}
            \forall j\in F_i, 
            w'(j) &=\bar w_i(j)
                &=(1-\alpha_i)\,w_i(j)
                &=r_i\,w_i(j)
        \end{aligned}
        \]
        (C2)が示された。
    \end{proof}

    \begin{body}
        \indent これが古典的チェバの定理の一般化になっていることは、前提・仮定・結論に含まれる命題の強弱を比較すれば直ちにわかる。

        \indent mathlib4内のCeva.Leanでは上記の主張をベースにLean4向けに整形した証明がなされている。

        \indent この定理を参考に、本稿では古典的チェバの定理の逆を一般化して示し形式化する。
    \end{body}

    % =========================================================
    \chapter{Lean4での実装}
    % =========================================================



    \section{チェバ型定理逆方向の主張と自然言語証明}
    \begin{body}
        \indent 古典的チェバの定理の逆を一般化すると、既存の順方向定理に対応する自然な主張は「ある大域的な重み $w$ が各局所的な重み $w_i$ と頂点 $i$ を除いて比例しているならば、そのチェビアンたちは共点するか、さもなくば互いに平行である」というものである。

        \indent 以下ではその主張を、順方向と同程度に一般的な形で述べる。
    \end{body}

    \begin{theorem}[チェバ型定理逆方向]
        \noindent\textbf{前提：}\par
        \begin{theoremenumerate}[label=(P\arabic*), leftmargin=4em, topsep=2pt, partopsep=0pt, itemsep=0pt, parsep=0pt]
            \item $k$ は体である。
            \item $V$ は $k$-ベクトル空間である。
            \item $P$ は $V$ 上のアフィン空間である。
            \item $\iota$ は有限の添字集合である。
            \item $s \subseteq \iota$ は空でない。
            \item 点族 $p:\iota\to P$ はアフィン独立である。
            \item 各 $i\in s$ について、重み $w_i:\iota\to k$ が与えられ、次を満たす：
            \[
            \sum_{j\in\iota} w_i(j)=1,
            \qquad
            w_i(i)=0,
            \qquad
            q(i)=\sum_{j\in\iota} w_i(j)\,p(j).
            \]
        \end{theoremenumerate}

        \noindent\textbf{仮定：}\par
        \begin{theoremenumerate}[label=(A\arabic*), leftmargin=4em, topsep=2pt, partopsep=0pt, itemsep=0pt, parsep=0pt]
            \item ある写像 $w:\iota\to k$ が存在して、次の条件を満たす：
            \[
            \forall i\in s,\ \exists r_i\in k,\ \forall j\in\iota,\ \bigl(j\neq i \Rightarrow w(j)=r_i\,w_i(j)\bigr).
            \]
            \item $w$ は零写像ではない：
            \[
            \exists j\in\iota,\ w(j)\neq 0.
            \]
        \end{theoremenumerate}

        \noindent\textbf{結論：}\par
        \begin{theoremenumerate}[label=(C\arabic*), leftmargin=4em, topsep=2pt, partopsep=0pt, itemsep=0pt, parsep=0pt]
            \item $\sum_{j\in\iota} w(j)\neq 0$ のとき、ある点 $p'\in P$ が存在して、任意の $i\in s$ について $p'$ は 2 点 $p(i)$ と $q(i)$ を結ぶ直線上に存在する。
            \item $\sum_{j\in\iota} w(j)=0$ のとき、任意の $i\in s$ について直線 $p(i)q(i)$ は互いに平行である。
        \end{theoremenumerate}
    \end{theorem}

    \begin{proof}
        $S:=\sum_{j\in\iota}w(j)$ とおく。

        まず各 $i\in s$ について、$w_i(i)=0$ および $\sum_j w_i(j)=1$ から
        \[
          q(i)-p(i)=\sum_{j\in\iota} w_i(j)\,(p(j)-p(i))
        \]
        が成り立つ。これは $q(i)$ を、基点を $p(i)$ とするアフィン結合として書き直した式である。

        \medskip
        \noindent\textbf{(1) $S\neq 0$ の場合}\par
        このとき正規化した重み $\widehat w:\iota\to k$ を
        \[
          \widehat w(j):=\frac{w(j)}{S}
        \]
        で定めると、明らかに $\sum_j \widehat w(j)=1$ である。よって
        \[
          p':=\sum_{j\in\iota}\widehat w(j)\,p(j)
        \]
        は点族 $p$ のアフィン結合として定まる。

        任意の $i\in s$ をとる。仮定 (A1) により、ある $r_i\in k$ が存在して $j\neq i$ ならば $w(j)=r_i w_i(j)$ である。そこで
        \[
        \begin{aligned}
          p'-p(i)
            &= \sum_{j\in\iota}\widehat w(j)\,(p(j)-p(i)) \\
            &= \frac{1}{S}\sum_{j\in\iota}w(j)\,(p(j)-p(i)) \\
            &= \frac{1}{S}\sum_{j\neq i}w(j)\,(p(j)-p(i)) \\
            &= \frac{r_i}{S}\sum_{j\neq i}w_i(j)\,(p(j)-p(i)) \\
            &= \frac{r_i}{S}\,(q(i)-p(i)).
        \end{aligned}
        \]
        したがって
        \[
          p' = \frac{r_i}{S}\,(q(i)-p(i)) + p(i)
        \]
        であり、$p'$ は $p(i)$ と $q(i)$ を通る直線上にある。$i\in s$ は任意であったから、$p'$ は全てのチェビアン $p(i)q(i)$ の共点である。

        \medskip
        \noindent\textbf{(2) $S=0$ の場合}\par
        $s$ は空でないから、ある $i_0\in s$ を固定する。ベクトル
        \[
          v:=\sum_{j\in\iota} w(j)\,(p(j)-p(i_0))\in V
        \]
        を考える。任意の $i\in s$ に対し、$S=0$ を用いると
        \[
        \begin{aligned}
          \sum_{j\in\iota} w(j)\,(p(j)-p(i))
            &= \sum_{j\in\iota} w(j)\,\bigl((p(j)-p(i_0))+(p(i_0)-p(i))\bigr) \\
            &= \sum_{j\in\iota} w(j)\,(p(j)-p(i_0)) + \left(\sum_{j\in\iota} w(j)\right)(p(i_0)-p(i)) \\
            &= v
        \end{aligned}
        \]
        となる。したがって各 $i\in s$ について
        \[
        \begin{aligned}
          v
            &= \sum_{j\in\iota} w(j)\,(p(j)-p(i)) \\
            &= \sum_{j\neq i} w(j)\,(p(j)-p(i)) \\
            &= r_i \sum_{j\neq i} w_i(j)\,(p(j)-p(i)) \\
            &= r_i\,(q(i)-p(i))
        \end{aligned}
        \]
        が成り立つ。よって各直線 $p(i)q(i)$ の方向ベクトル $q(i)-p(i)$ は、同じベクトル $v$ のスカラー倍である。

        さらに $v\neq 0$ である。実際、もし $v=0$ ならば
        \[
          \sum_{j\in\iota} w(j)=0,
          \qquad
          \sum_{j\in\iota} w(j)\,(p(j)-p(i_0))=0
        \]
        となり、第 $4$ 章で述べたアフィン独立性の系より全ての $j\in\iota$ で $w(j)=0$ となる。これは (A2) に反する。

        したがって $v$ は非零であり、全ての直線 $p(i)q(i)$ は同じ非零方向ベクトルに平行である。よってこれらの直線は互いに平行である。
    \end{proof}

    \begin{body}
        \indent この定理では、大域的な重み $w$ の総和 $\sum_j w(j)$ が 0 かどうかが「有限の共点」と「無限遠方向」を分けている。総和が非零なら重みを正規化して実際の点 $p'$ を作れる。一方、総和が 0 だと点そのものは得られないが、零和の重み付き差分が共通方向ベクトルを与えるため、チェビアンは互いに平行になる。
    \end{body}

    \section{設計}
    \subsection{設計の問題点}
    \begin{body}
        \indent 素朴には上記の定理をそのままLean4に書き直せば良いと考えられる。しかしここで生じる問題は、$4$章で行った概念の定義は本稿で証明を行いやすくするために採用したものであり、mathlib4中に用意された概念の定義とはいささか異なる。例えば、チェバ型定理という単一の定理を紙面上で示すと言う目的でアフィン結合$\to$アフィン部分空間という定義の流れを本稿では選択したが、これは$V$が$k$-加群であることを先に仮定しているという意味で数学的にあまり綺麗な方法ではない。実際mathlib4中の定義ではアフィン部分空間$\to$アフィン結合の流れが採用されており、本稿で用いた定義とは流れも見た目も異なっている概念も多い。よって、形式化の際はこのような定義の読み替えを適宜行いつつ翻訳を行っていく必要がある。
    \end{body}

    \subsection{設計方針と採用理由}
    \begin{body}
        \indent 前章で見たようにJoseph Myers氏は最も一般的な定理から徐々に特殊な場合の定理を導くという流れをとっており、チェバ型定理逆方向の実装でも同様の方法を取ることにする。

        \indent 具体的には、以下の通りである。

        \begin{bodyenumerate}
            \item (2)を示すための補題
            \item 任意添字点族チェバの定理の逆
            \item 有限添字点族チェバの定理の逆
            \item 有限添字点族チェバの定理の逆（平行を除いた版）
            \item $2$-単体についての古典的チェバの定理の逆
            \item $2$-単体についての古典的チェバの定理の逆（平行を除いた版）
        \end{bodyenumerate}

        % \indent 基本的に構成は同じであり、

    \end{body}


    \section{追加モジュールの実装}

    \subsection{概念と既存との対応}
    \begin{body}
        \indent 第 $4$ 章で紹介した数学的概念や、その証明で実際に使う補助道具が mathlib4 上でどのように実装されているかを整理する。
        \begin{bodydescription}[style=nextline,leftmargin=4.5em,labelsep=1em,font=\normalfont\bfseries,itemsep=0pt, topsep=2pt, parsep=0pt]
          \item[アフィン空間] \texttt{AddTorsor V P}, \texttt{AffineSpace V P}
          \item[アフィン独立] \texttt{AffineIndependent k p}
          \item[アフィン結合] \texttt{Finset.affineCombination k s p w}
          \item[零和の重み付き差] \nolinkurl{Finset.weightedVSub}, \nolinkurl{Finset.weightedVSubOfPoint}
          \item[重心座標] \texttt{AffineBasis.coord}
          \item[アフィン包] \texttt{affineSpan k (s : Set P)}
          \item[直線上の点の表示] \nolinkurl{AffineMap.lineMap}, \nolinkurl{mem_affineSpan_pair_iff_exists_lineMap_eq}, \nolinkurl{mem_affineSpan_pair_iff_exists_lineMap_rev_eq}
          \item[平行] \texttt{AffineSubspace.Parallel}, 記法 $\parallel$
          \item[凸包] 一般の凸包は \texttt{convexHull k (s : Set V)} で与えられる。単体の頂点が張る凸包に対応するものとしては \texttt{Affine.Simplex.closedInterior} が用意されている。
          \item[単体] \texttt{Affine.Simplex k P n}
          \item[面] \texttt{Affine.Simplex.face}
          \item[対向面] \texttt{Affine.Simplex.faceOpposite}
        \end{bodydescription}

        \indent とくに逆方向定理の証明で重要なのは、(i) アフィン結合を正規化して点を作る部分、(ii) 零和の重み付き差を共通方向ベクトルとして扱う部分、(iii) 「ある点が 2 点を結ぶ直線上にある」ことを \nolinkurl{AffineMap.lineMap} や 2 点のアフィン包への所属に落とす部分、の 3 箇所である。
    \end{body}


    \subsection{追加するモジュール}
    \begin{body}
        \indent 既存の mathlib4 だけでも順方向チェバの定理を使うための土台は整っているが、逆方向を読みやすく実装するにはいくつかの補助補題を追加した方がよい。重要なのは新しいデータ型を導入することではなく、既存 API の上に薄いラッパーを重ねることである。
    \end{body}

    \begin{bodydescription}
        \item[\hspace*{1.5em}・重みの正規化に関する補題]
        $S=\sum_j w(j)\neq 0$ のとき、$w/S$ が和 $1$ を持つこと、また正規化後の重みが各 $w_i$ と頂点 $i$ を除いて比例し続けることをまとめて扱う補題を用意する。これは共点 $p'$ を \nolinkurl{Finset.affineCombination} で構成する箇所を簡潔にする。

        \item[\hspace*{1.5em}・零和重みから共通方向を作る補題]
        $\sum_j w(j)=0$ のとき、\nolinkurl{Finset.weightedVSub} が基点の選び方に依らないことを使って共通方向ベクトルを定義する補題を用意する。さらに、アフィン独立性からそのベクトルが $0$ でないことを導く補題が必要である。これが「平行」ケースの核心になる。

        \item[\hspace*{1.5em}・直線所属と \nolinkurl{lineMap} の往復補題]
        自然言語では「$p'$ は $p(i)$ と $q(i)$ を結ぶ直線上にある」と書けば十分だが、Lean ではこれを 2 点のアフィン包への所属あるいは \nolinkurl{AffineMap.lineMap} による表示へ落とした方が扱いやすい。既存の \nolinkurl{mem_affineSpan_pair_iff_exists_lineMap_eq} を、今回の証明で直接使える形に整える小補題を追加する。

        \item[\hspace*{1.5em}・特殊化のための橋渡し補題]
        逆方向も順方向と同様に、一般添字版 $\to$ 有限添字版 $\to$ 三角形版、という流れで特殊化できるようにしておくと再利用性が高い。したがって本体の定理のほかに、有限型への読み替えと三角形版への特殊化を担う橋渡し補題を用意するのが望ましい。
    \end{bodydescription}

    \begin{body}
        \indent 実装上は、これらを 1 つの大きなファイルで一気に証明するよりも、既存の \nolinkurl{Combination}・\nolinkurl{AffineSubspace.Basic}・\nolinkurl{Ceva} を import した補助ファイルに分割し、レビュー可能な小さな PR 単位に切り出していく方が現実的である。逆方向そのものは新しい理論分野を作るものではなく、既存のアフィン幾何 API の「つなぎ方」を整える作業として理解するのが適切である。
    \end{body}

    % \subsection{実装の工夫事項・特記事項}
    % \begin{bodydescription}
    %     \item \textbf{実装を急いでいます}
    % \end{bodydescription}

    % =========================================================
    \chapter{結論と今後の展望}
    % =========================================================


    \section{結論}
    \begin{body}
        \indent 本稿では、既に mathlib4 に導入された一般化チェバの定理の順方向を整理した上で、その逆方向として自然な「共点または平行」版を定式化した。さらに、その自然言語証明を与え、mathlib4 上で実装する際に必要となる既存概念との対応、および補助補題の設計方針を明確にした。
    \end{body}

    \section{今後の課題}
    \begin{body}
        \indent 当初は一般化チェバの定理の順方向も含めて実装する予定だったが、コミュニケーション不足により Joseph Myers 氏の一般化順方向 PR が先行して導入された。結果として本稿では逆方向に焦点を絞ることになり、研究内容を早い段階で共有することの重要性を痛感した。

        \indent また、数学にしろ情報学にしろ、何を達成したいかを正確に把握して最適な言語・理論を選択することが不可欠であることを痛感した。単一の定理の自然言語証明を組み立てるための定義と、再利用可能なライブラリ API として採用すべき定義は必ずしも一致しない。自然言語・一階述語論理・型理論のいずれで記述するかによっても見通しの良い定式化は変わるため、その翻訳方針の選定が本研究で最も難しい点であった。
        
        \indent 今回培ったOSS開発の作法に従い、次回以降のプロジェクトにも取り組んでいきたい。
    \end{body}

    \section{Lean4の普及と形式化研究の展望}
    \begin{body}
        \indent 現在LeanのAutoformalizationという技術開発が進んでおり、自然言語で記述した数学の証明をそのままLeanの形式言語に書き換えることが目指されている。\\
        今回は個別の定理を形式化するというプロジェクトであったが、この技術を用いればある数学的な証明を自然言語で記述するだけで、立ちどころに形式証明が得られると考えられている。
        この技術は、Lean研究を飛躍的に促進すると考えられる。個別の定理の形式化だけでなく、このAutoformalizationに関する見識も深めていきたい。
    \end{body}

    % =========================================================
    \appendix
    % =========================================================

    % =========================================================
    % 付録：型理論についての補足
    % =========================================================
    \chapter{型理論についての補足}\label{appendix:type-theory}

    \section{本付録の目的}
    \begin{body}
        \indent 本章で触れたように型理論では、命題を型・証明をその型の項、として表現する。しかしこれだけでは主に一階述語論理ベースの集合論に慣れている読者にはその全貌は見せられないように思う。よってよりイメージのつきやすい型理論の説明を一階述語論理と比較し、形式体系とは何かに立ち返りながら行うことが本付録の目的である。
    \end{body}

    \section{形式体系とは}

    \begin{body}
        \indent 数学の推論を厳密に形式化して分析する分野を数理論理学と呼び、形式体系とはその中で用いる概念の一つである。数理論理学は構文論と意味論に分かれるが、両者は形式体系を別のアプローチで分析する。構文論とは記号列の操作（導出）の規則に着目する手法であり、意味論とは記号列に解釈を与えてその真偽を論じる手法である。

        \indent 一階述語論理や高階述語論理、型理論など、形式体系と呼ばれるものは様々あり、形式言語を用いて数学の証明・推論を様々な方法で表現する。構文論的な観点で見たときこれらの多くに共通する眼目は、証明という対象を有限のデータとして扱う点にある、つまり\emph{証明とは有限の記号列でなければならない}。
    
        \indent 形式体系はしばしば \emph{言語}、\emph{公理}、\emph{推論規則}の三つ組として記述される。言語は「許される記号と、それらから整式（文法を満たす記号列）を作る規則」を与える。公理は「証明の出発点として採用する式またはその集合」である。推論規則は「前提となる式から結論となる式を得る規則」である。これら自体は有限である必要はない。実際、変数記号は自然数の個数だけ無数に用意できるし、公理はスキームという形で実質無限本存在する、その上推論規則の中には$\omega$規則$\frac{\{\varphi(n)\}_{n\in\mathbb{N}}}{\forall n\in\mathbb{N}\,\varphi(n)}\;(\omega)$と呼ばれるものもあり、推論が無限を含んでいる。

        \indent しかし証明を構築するときは話が別である。証明とは「人間が数学的主張の正しさを検証するための道具」であるから有限でなくては検証ができない。よっていざ証明を構築するときは、言語から有限個の記号を選び、公理スキームの有限個を利用し、有限個の仮定から有限個の結論が得られる推論規則のみを利用する。したがって形式体系を構築するにあたって、言語の無限個の記号、公理の無限本の公理図式は許容されるが、$\omega$規則のような\emph{無限推論規則}は「証明が有限」という前提自体を捨てる必要が生まれるため一般的な形式体系の議論では許容されない。（これを認めた無限論理・無限証明論と呼ばれる分野も存在するが、無限を扱える代わりに完全性やコンパクト性などの豊かな性質の多くが失われるまたは変質してしまう別物として区別され研究されている）
    
        \indent 本稿では、与えられた公理（集合）に推論規則を適用して得られる導出可能な式全体の閉包を \emph{理論体系}と呼ぶ。「これらの式が何を意味するか」という解釈を与えるのは意味論に譲る。
    
        \indent 例えば、多くの数学書が暗黙のうちに議論に用いている集合論（典型的には ZF / ZFC）は、形式体系としては「集合 $\in$ を含む言語」「公理（外延性・分出・置換スキーム等）」「古典論理の推論規則」を備えた理論として以下のように整理できる。
        
        \vskip\baselineskip

        \noindent\textbf{集合論}
        \begin{bodyitemize}[leftmargin=7em]
            \item[言\hspace*{2em}語] : 一階述語論理記号$\cup$変数記号$\cup$補助記号$\cup\in$と論理式の定義
            \item[推論規則] : 自然演繹
            \item[公\hspace*{2em}理] : ZF / ZFC
        \end{bodyitemize}

        \vskip\baselineskip

        \indent まず、言語を定義する。記号集合$\{\neg,\to,\forall,\exists$,…$\}\cup \{v_{0},v_{1},v_{2}$,…$\}\cup \{\in\}$を用意し、この集合から$\neg v_{2}\forall, \forall v_{0}(v_{0}=v_{0}), \cup v_{n}\cap v_{n-1}=v_{0}$のように無数の記号列を構成できる。そして得られた記号列から記号列として認めるものを定義する。認めた記号列を論理式と呼び$\phi,\psi,\varphi$などと略記する。先ほどの例では$\forall x(x=x)$のみ論理式である。\\
        次に、推論規則を定義する。有名な例では含意除去$\infer[\to E]{\varphi \\ \varphi \to \psi}{\psi}$（Modus Ponensのこと）などがある。先ほどの例に対しては例えば全称除去$\infer[\forall E]{\forall x\,\varphi(x)}{\varphi(t)}$を用いて$v_{0}=v_{0}$などの新たな論理式が得られる。\\
        最後に公理を決める。Kunenの定義に従えば、外延性の公理
        \[
        \forall v_0\,\forall v_1\;\Bigl(\forall v_2\,(v_2\in v_0 \leftrightarrow v_2\in v_1)\;\rightarrow\; v_0=v_1\Bigr)
        \]
        などがある。論理式の中で証明の出発点として採用する特別な論理式である。

        \indent こうして得られた公理から推論規則による変形を通して得られたものだけが定理と呼ばれ、理論体系を構成する一員となる。

        \indent こう書くと、自然に「では言語、公理、推論規則のいずれかを変更したらどうなるだろう？」という疑問が自然に湧いてくるであろう。公理的集合論では主にこの公理の部分を様々に入れ替え、付け加え、除くことで理論体系がどう変化するかを分析する。また、推論規則としてはヒルベルト式やシーケント計算が存在し定理の証明のしやすさに影響を与える。

        \indent 同様にしてこの三つ組を変更して得られるのが、型理論である。（正確にいうと、上記の三つ組は一階述語論理ベースの形式体系を分類するには有用だが、型理論ベースの形式体系を分類するには不向きである。その上で敢えて同じ視点でCICを眺めてみる。）


        以下では、集合論に慣れた読者向けに、CIC を「判断を生成する規則の体系」として見たときの要点を順に説明する。
        \vskip\baselineskip
        \begin{bodyitemize}[leftmargin=4em]
            \item[\textbf{CIC}] \mbox{}
            \item[言語]
                文脈$\Gamma$・項$t$・型$A$からなる依存型付き$\lambda$計算の構文
            \item[公理]
                特になし
                ただし必要に応じて定数宣言$c:A$を追加して公理的拡張を行うことがある。
            \item[推論規則]
                判断
                \[
                \Gamma \ \mathsf{ctx},\qquad
                \Gamma \vdash A \ \mathsf{type},\qquad
                \Gamma \vdash t : A
                \]
                を生成する型付け規則（変数規則・$\Pi$の形成/導入/除去・帰納型の形成/導入/除去など）、計算規則（$\beta$簡約・帰納型の$\iota$計算規則等）および、変換規則
        \end{bodyitemize}

        \indent 見慣れない用語が並びイメージが付きづらいように思う。一階述語論理に慣れている読者が型理論をゼロから捉えるうえで重要なのは、集合論のように「ある宇宙をまず仮定して、その上で式の真偽を語る」という発想をいったん脇に置くことである。型理論（少なくともCICの系譜）では、最初にあるのは\emph{対象世界}ではなく、\emph{判断を生成する推論体系}である。つまり「何が式（項・型）として許されるか」「その式がどの型を持つと主張してよいか」「二つの式が同一視されてよいか」を、構文と規則によって定義する。集合論が「$\in$ を含む一階言語 + 公理 + 推論規則」で理論を作るのと同様に、型理論は「判断（judgement）+ 生成規則」で理論を作る。

        \indent まず\emph{判断}の形だけを固定する。判断とはメタ主張のことを指す。CICの核は典型的に次の三種の判断で動く：
        (1) 文脈が整っている $\Gamma\ \mathsf{ctx}$、
        (2) 文脈の下で型が整っている $\Gamma\vdash A\ \mathsf{type}$、
        (3) 文脈の下で項が型を持つ $\Gamma\vdash t:A$。
        ここで「文脈」は自由変数とその型の並び（および場合によっては定義）であり、一階論理の「変数割当」や「仮定集合」に対応する。だが対応は完全ではない。型理論では文脈そのものが推論の対象であり、文脈拡張（$\Gamma,x:A$）が推論規則として明示される。

        \indent 次に、\emph{構文の生成}を与える。依存型付き$\lambda$計算では、項は変数、適用 $t\,u$、抽象 $\lambda x:A.\,t$ 等から作られ、型は関数型（依存積）$\Pi x:A.\,B$ 等から作られる。ここで本質は「型が項に依存できる」ことで、$\Pi x:A.\,B$ に現れる $B$ は $x$ を含みうる。これが、一階論理で言えば「述語が変数に依存する」ことに対応するが、型理論では依存が\emph{型形成のレベル}で起きる点が決定的に違う。

        \indent しかし構文だけでは何も主張できない。そこで三段階目として、\emph{型付け規則（判断生成規則）}を置く。最も基本的なのは変数規則で、$\Gamma$ が整っていて $\Gamma$ に $x:A$ が含まれるなら $\Gamma\vdash x:A$ を得る。続いて、$\Pi$型について形成・導入・除去（いわゆる関数型の形成、ラムダ抽象、適用）が与えられる。これらは自然演繹の対応物として読むのが早い：$\Pi$は全称量化と含意を統合したような振る舞いをし、$\lambda$は仮定導入、適用は仮定消去に相当する。ただし「論理式の証明」という読みは\emph{後から可能になる解釈}であって、体系の定義そのものはあくまで型付け規則の集合として完結する。

        \indent 四段階目が、集合論側の読者が最初に転ぶ地点である。型理論には「等号」が二種類ある。ひとつは\emph{定義的等しさ}（definitional equality）で、計算により同一視してよいという約束である。$\beta$簡約（$(\lambda x.\,t)\,u \leadsto t[x:=u]$）や帰納型の計算規則（$\iota$規則）がここに入る。もうひとつは\emph{命題としての等しさ}（identity type）で、$t=u$ を証明（項）として与える対象である。CICの核（カーネル）が自動で判定するのは前者であり、後者は通常の命題と同様に証明が必要になる。ここを混同すると、型理論が「計算の体系」なのか「論理の体系」なのか見失う。実態は両方で、計算が\emph{型付けの同値判定}に組み込まれている点が特徴である。

        \indent 五段階目として、\emph{帰納型}を導入する。CICが「表現力の強い数学の言語」たりうるのは、自然数・リスト・木・有限集合・（適切な形の）集合様構造などを、帰納型の形成規則とその帰納法/再帰原理として内蔵するからである。ここで公理を足すのではなく、\emph{新しい型構成子とその推論規則}を足す、という設計思想が見える。集合論で「分出・置換・無限」等を公理として積むのに対して、型理論では「このデータ型をこういう生成子で作る」「それに対する関数は再帰で定義できる」という規則として積む。

        \indent 六段階目が\emph{宇宙}である。素朴に「型の全体」を一つの型として置くとラッセル型の破綻が起きるため、CICでは $Type_0 : Type_1 : Type_2 : \cdots$ のような階層を置き、型がどの宇宙に属するかを追跡する。これは集合論の累積階層 $V_\alpha$ を思い出せば直観は掴めるが、同一ではない。宇宙は「型の型」というメタ構造を\emph{対象言語内に反映するための装置}であり、同時に矛盾回避のための層構造である。

        \indent ここまでで、公理が「特になし」と書かれている意味が見えてくる。CICは、何らかの特定の数学的対象（例えば $\mathbb{N}$ の標準モデル）を最初に仮定しない。代わりに、判断生成規則が許す限りの構成を行う。その上で、必要なら定数宣言 $c:A$ を追加して拡張する。これは集合論で言えば「新しい定数記号を導入し、追加公理でそれを規定する」操作に似ている。ただしCICでは、その追加は「証明できないが仮定したいもの（例えば選択公理、古典命題の排中律、特定の外延性原理など）」を\emph{公理として型に住まわせる}形で行われる。どの公理を足すかで、計算性・構成性・抽出可能性といった性質が大きく変わるので、無自覚に足すと設計が壊れる。

        \indent 以上を踏まえると、CICは「一階論理 + $\in$」のような単一の基礎記号から世界を立てる流儀ではなく、「判断の生成規則（型付け）+ 計算（簡約）+ 帰納原理」という三点セットから世界を立てる流儀だと言える。集合論が\emph{真理値}を中心に体系を組むのに対し、型理論は\emph{構成（項）}を中心に体系を組む。命題＝型、証明＝項という対応はこの上に乗る見方であり、最初に置くべき土台は「判断生成規則の体系」である。

        \indent 型理論は一階述語論理とは異なる観点から定義される形式体系を提供し、数学の証明とは何かを探る重要なツールとなっている。
    \end{body}

    % =========================================================
    % 付録：　古典的チェバの定理の版の教科書による違いについて
    % =========================================================
    \chapter{古典的チェバの定理の版の教科書による違いについて}\label{appendix:classical-ceva-version}

    \section{本付録の目的}
    \begin{body}
        \indent 本付録では、高校数学で学ぶ古典的チェバの定理の版の違いがなぜ生まれ、どのような意図があるのかを解説する。
    \end{body}

    \section{古典的チェバの定理の版の種類}
    \begin{body}
        \indent 古典的チェバの定理は大きく分けて以下の3つの形で紹介されうる。
        \begin{bodyitemize}
            \item[①] $D, E, F$を辺上に限定した、3直線が1点で交わるための必要十分条件
            \item[②] $D, E, F$を辺の延長上でも許した、3直線が1点で交わる必要条件(\hyperref[thm:classical-ceva]{定理3.1})
            \item[③] $D, E, F$を辺の延長上でも許した、3直線が1点で交わる必要十分条件
        \end{bodyitemize}

        \indent 一見それぞれ異なる主張をしているように見えるがこれらは数学的に何も問題はなく、\emph{背後の幾何学体系に何を採用するか}で主張の範囲や述べ方が違うだけである。

    \end{body}

    \section{古典的チェバに版の違いが生まれる要因}
    \begin{body}
        \indent ①〜③の中で本稿で取り扱っている②のみ必要十分性を主張する定理になっていない。試しに定理の安直な逆命題を考えてみると成立しないことがわかる。その原因となっているのが、下図のような、チェビアンが平行になる場合である。
    \end{body}

    \begin{center}
        \begin{tikzpicture}[scale=0.8, x=1cm,y=1cm, line cap=round, line join=round, >=Latex]
            % ---- keep the whole picture nicely centered ----
            \path[use as bounding box] (-1.2,-1) rectangle (7.2,10.2);
    
            % --- Title (Added per request) ---
            \node[font=\bfseries] at (3,9.8) {(図)等式条件を満たしながらチェビアンが平行となる場合};
    
            % --- coordinates ---
            \coordinate (A) at (0,0);
            \coordinate (B) at (6,0);
            \coordinate (C) at (4,3);
    
            \coordinate (D) at (0,9);
            \coordinate (E) at (6,4.5);
            \coordinate (F) at (4,0);
    
            % --- triangle ---
            \draw[thick] (A) -- (B) -- (C) -- cycle;
    
            % --- side extensions (dashed) ---
            \draw[dashed] (B) -- (D);  % extension of BC to meet x=0
            \draw[dashed] (A) -- (E);  % extension of CA to meet x=6
    
            % --- points ---
            \fill (A) circle (1.6pt) node[below left] {$A$};
            \fill (B) circle (1.6pt) node[below right] {$B$};
            \fill (C) circle (1.6pt) node[above, xshift=2pt] {$C$};
    
            \fill (D) circle (1.6pt) node[left] {$D$};
            \fill (E) circle (1.6pt) node[right] {$E$};
            \fill (F) circle (1.6pt) node[below] {$F$};
    
            % --- Improved Parallel Mark Style ---
            \tikzset{
              parallelArrow/.style={
                postaction={decorate},
                decoration={
                  markings,
                  mark=at position 0.5 with {
                    \draw[line width=0.8pt] (-3pt,-3pt) -- (0pt,0pt) -- (-3pt,3pt);
                  }
                }
              }
            }
    
            % --- cevians (parallel) with fixed arrows ---
            \draw[very thick, parallelArrow] (A) -- (D);
            \draw[very thick, parallelArrow] (B) -- (E);
            \draw[very thick, parallelArrow] (F) -- (C); 
    
        \end{tikzpicture}
    \end{center}

    \begin{body}
        \indent この場合、等式条件は成立するがチェビアンが平行となってしまい共点していない。まさにこの問題によって、背景とする幾何学体系による定理の述べ方が変わるのである。
    \end{body}

    \section{各々の版の背景}

    \begin{body}
        \indent 結論から述べると、①〜③はそれぞれユークリッド幾何、アフィン幾何、射影幾何の観点から古典的チェバの定理を述べたものである。

        \indent ①はユークリッド幾何の観点で述べた形である、$BD$、$DC$などは全てユークリッド距離であり常に正となる。ユークリッド距離ではチェビアンの足が対辺を外分するときをうまく表現できないため、あらかじめチェビアンの足は対辺中に限定される。チェビアンが平行になることは排除できないが、その場合等式条件が成立しないため必要十分性を主張する定理となる。

        \indent ②はアフィン幾何の観点で述べた形である。ここで登場するのはアフィン比と呼ばれる$\dfrac{BD}{DC}$という値であり、辺を成す$2$頂点の重心座標から得られるものである。高校範囲内ではアフィン幾何を登場させられないため、ユークリッド幾何の範疇で定義できる有向線分という考え方を用いて表現している。しかしこの場合定理の逆がそのままでは成り立たず、等式条件が成立してもチェビアンが平行になってしまうという反例が生じる。

        \indent ③は射影幾何の観点で述べた形である。射影幾何はユークリッド幾何やアフィン幾何を包含する幾何学の大きな枠組みの一つである。重要な特徴が空間内に無限遠点を含むため、平行という概念を「無限遠点で交点を持つ」と言える点である。定理の主張自体は②と変わっていないが、交点をもつという言葉の意味が広がったことで述べ方が変わっている。

        \indent 高校数学で習うのは高々ユークリッド幾何であり違いを説明するのが困難なため、追加の前提条件を加えるなどして主張が偽にならないように工夫して解説するという教育的な配慮がなされている。

        \indent 本稿ではアフィン幾何に基づいて議論を進めた。「チェビアンが平行でない」という趣旨の前提を加えれば強引に必要十分性を主張する定理にできるが、定理の自然さが損なわれると考え②の形を選択した。
    \end{body}

    \backmatter
    \addchaptertotoc{謝辞}
    \begin{body}
        \indent 本研究を進めるにあたり、終始丁寧なご指導とご助言を賜りました薄葉季路教授に心より感謝申し上げます。研究の方向性に迷った際にも温かく見守りながら的確な示唆をくださり、本研究を完成させることができました。

        \indent また、Lean4およびmathlib4のコミュニティの皆様には、PRのレビューや議論を通じて多くの学びをいただきました。特にJoseph Myers氏には、一般化チェバの定理の形式化に関する先行的な取り組みを通じて大きな刺激を受けました。

        \indent 最後に、日頃より支えてくださった家族と友人に感謝いたします。
    \end{body}


    \begin{thebibliography}{99}
        \bibitem{Kunen} Kenneth Kunen, \emph{Set Theory: An Introduction to Independence Proofs}, North-Holland, 1980.
        \bibitem{LeanProverCommunity} Lean Prover Community, \emph{Lean Prover Community}, \url{https://leanprover-community.github.io/}, accessed 2026-01-18.
        \bibitem{DeepMindIMO2025} Google DeepMind, \emph{Advanced version of Gemini with Deep Think officially achieves gold-medal standard at the International Mathematical Olympiad}, 2025.
        \bibitem{multipede} Dov Samet, \emph{An Extension of Ceva's Theorem to $n$-Simplices}, 2021.
        \bibitem{LeanResearch} 秋田 隼, \emph{LeanResearch}, GitHub repository, \url{https://github.com/Aj1905/LeanResearch.git}, 2026.
        \bibitem{mathlib3_ceva} Mantas Bak\v{s}ys, \emph{Ceva's Theorem (mathlib3)}, GitHub Pull Request \#10632, \url{https://github.com/leanprover-community/mathlib3/pull/10632}, 2021--2023.
        \bibitem{mathlib4_ceva_docs} Lean Prover Community, \emph{Mathlib.LinearAlgebra.AffineSpace.Ceva}, mathlib4 documentation, \url{https://leanprover-community.github.io/mathlib4_docs/Mathlib/LinearAlgebra/AffineSpace/Ceva.html}, accessed 2026-01-18.
        \bibitem{mathlib4_ceva_myers} Joseph Myers, \emph{feat(LinearAlgebra/AffineSpace/Ceva): Ceva's theorem}, GitHub Pull Request \#33409, \url{https://github.com/leanprover-community/mathlib4/pull/33409}, merged 2026-01-14.
    \end{thebibliography}

\end{document}
\end{document}
